% !PRINTED 105
\weberpage{105}
%<<CONT
\weq{11}{\begin{aligned}
e^{-\frac{\pi i u^2}{\omega}}\vartheta_{11}\left(\frac{u}{\omega}, -\frac{1}{\omega}\right) &= -i\sqrt{-i\omega}\,\vartheta_{11}(u), \\
e^{-\frac{\pi i u^2}{\omega}}\vartheta_{01}\left(\frac{u}{\omega}, -\frac{1}{\omega}\right) &= \sqrt{-i\omega}\,\vartheta_{10}(u), \\
e^{-\frac{\pi i u^2}{\omega}}\vartheta_{10}\left(\frac{u}{\omega}, -\frac{1}{\omega}\right) &= \sqrt{-i\omega}\,\vartheta_{01}(u), \\
e^{-\frac{\pi i u^2}{\omega}}\vartheta_{00}\left(\frac{u}{\omega}, -\frac{1}{\omega}\right) &= \sqrt{-i\omega}\,\vartheta_{00}(u).
\end{aligned}}

\begin{center}
\textbf{§ 32. \gesperrt{Die Haupttransformationen zweiter Ordnung der $\vartheta$-Funktionen.}}
\end{center}

Die beiden Haupttransformationen $n$ter Ordnung
\[
\begin{pmatrix} 1, & 0 \\ 0, & n \end{pmatrix}, \quad \begin{pmatrix} n, & 0 \\ 0, & 1 \end{pmatrix}
\]
verwandeln nach § 27, (11) die Charakteristik $(g_1, g_2)$ in
\[
(ng_1, g_2), \quad (g_1, ng_2),
\]
d. h. bei ungeradem $n$ bleibt die Charakteristik ungeändert, bei geradem $n$ geht sie über in
\weq{1}{(0, g_2) \text{ oder } (g_1, 0).}
Da sich hiernach die Transformation geraden Grades wesentlich anders verhält als die ungeraden Grades, so betrachten wir zunächst den Fall $n = 2$. Die erste und zweite Haupttransformation zweiten Grades werden die \gesperrt{Landensche} und die \gesperrt{Gausssche} Transformation genannt.

Nach § 27, (12) sind
\weq{2}{\begin{aligned}
\vartheta_{g_1, g_2}\left(u, \frac{\omega}{2}\right) &= \Theta_{g_1, 0}(u, \omega) \\
\vartheta_{g_1, g_2}(2u, 2\omega) &= \Theta_{0, g_2}(u, \omega)
\end{aligned}}
$\Theta$-Funktionen zweiter Ordnung von $u, \omega$, die sich nach § 21 darstellen lassen.

Wir erhalten zunächst die zwei Formelpaare, in denen $A, B$ von $u$ unabhängig sind:
\weq{3}{\begin{aligned}
A\vartheta_{11}\left(u, \frac{\omega}{2}\right) &= \vartheta_{01}(u, \omega)\vartheta_{11}(u, \omega), \\
A\vartheta_{10}\left(u, \frac{\omega}{2}\right) &= \vartheta_{00}(u, \omega)\vartheta_{10}(u, \omega),
\end{aligned}}
\weq{4}{\begin{aligned}
B\vartheta_{11}(2u, 2\omega) &= \vartheta_{10}(u, \omega)\vartheta_{11}(u, \omega), \\
B\vartheta_{01}(2u, 2\omega) &= \vartheta_{00}(u, \omega)\vartheta_{01}(u, \omega),
\end{aligned}}
%CONT>>
